% =============================================================================
%  SPECTRA — Section VI: Results and Discussion
%  Target: IEEE Transactions on Information Forensics & Security
%  Author: Sagar Korde
% =============================================================================
%
%  Required packages (in main .tex preamble):
%    \usepackage{booktabs}
%    \usepackage{siunitx}
%    \usepackage{multirow}
%    \usepackage{threeparttable}
%    \usepackage{xcolor}
%    \usepackage{amsmath}
%
%  \sisetup{round-mode=places, round-precision=4}
% =============================================================================

\section{Results and Discussion}
\label{sec:results}

We organize the evaluation around four axes:
(i)~tabular classification performance (Track~T1),
(ii)~graph-topology classification performance (Track~T2),
(iii)~temporally-fair ablation, and
(iv)~unsupervised clustering quality.
Section~\ref{subsec:spectra_outputs} then examines SPECTRA-specific
diagnostic outputs (drift, trust, Markov anomaly) that are qualitatively
distinct from any single-objective baseline.
All numbers reported here are computed on the held-out test set;
no hyperparameter selection was performed after observing test outcomes.

% ─── VI-A  Tabular Classification ────────────────────────────────────────────

\subsection{Classification Performance (Track T1: Tabular)}
\label{subsec:t1_results}

Table~\ref{tab:classification} (upper block) reports Track~T1 results for
the two tabular classifiers evaluated on the original \num{20}-feature set
$\mathbf{X}^{\dagger}_{\text{scaled}}$.
Both XGBoost and BERT4ETH attain perfect scores across every metric
($\text{Acc} = \text{F1} = \text{AUC-ROC} = \text{MCC} = \num{1.0000}$).

\textbf{Label-feature entanglement, not generalization.}
These oracle-level scores do \emph{not} represent genuine generalization.
As established formally in Section~V-B, the CryptoGraphNet labels are
deterministically derived from the same structural fields that form the
feature set (input/output counts, fee patterns, script type combinations).
A classifier trained on these features is therefore solving a
deterministic reconstruction problem rather than learning a probabilistic
decision boundary.
The confusion matrix is trivially zero off-diagonal: given any legal feature
vector, the label is algebraically recoverable.
We retain these results for completeness and to establish the hard upper
bound imposed by label-feature entanglement; they should not be interpreted
as evidence that these classifiers generalize to unseen transaction patterns
drawn from a different annotation schema.

\textbf{Computational cost.}
The two oracles incur radically different training costs for identical
predictive outcomes.
XGBoost converges in \num{6.70}\,s with a per-sample inference latency
of \num{0.0013}\,ms.
BERT4ETH requires \num{229.5}\,s of fine-tuning---a \num{34}$\times$ overhead---
and achieves \num{0.0017}\,ms inference latency.
The BERT4ETH transformer architecture provides no marginal accuracy benefit
over a gradient-boosted tree on this tabular task, a finding consistent with
recent benchmarks on structured tabular data~\cite{Grinsztajn2022}.

% ─── VI-B  Graph-Topology Classification ─────────────────────────────────────

\subsection{Classification Performance (Track T2: Graph-Topology)}
\label{subsec:t2_results}

Table~\ref{tab:classification} (lower block) reports Track~T2 results for
methods that operate on the address graph $\mathcal{G}$ and spectral
embeddings $\boldsymbol{\Phi}$ rather than on raw transaction tabular features.
This is the methodologically sound comparison regime: no method has access
to the label-entangled transaction fields.

\textbf{SPECTRA-full outperforms all graph baselines.}
The full SPECTRA ensemble achieves $\text{Acc} = \num{0.9160}$ and
$\text{MCC} = \num{0.8397}$, the highest values in the graph-topology track.
This represents a margin of $+\num{33.9}$ percentage points (pp) over
EvolveGCN ($\text{Acc} = \num{0.5773}$) and $+\num{41.7}$\,pp over
AA-GNN ($\text{Acc} = \num{0.4992}$).
In terms of discriminative power, MCC is the most informative scalar
summary under class imbalance~\cite{Chicco2020}; SPECTRA-full's MCC of
\num{0.8397} versus EvolveGCN's \num{0.5002} and AA-GNN's \num{0.4101}
confirms that the performance gap is not an artifact of majority-class
dominance.

\textbf{GNN module alone versus full ensemble.}
The SPECTRA GNN module operating in isolation achieves only
$\text{Acc} = \num{0.4036}$ and $\text{MCC} = \num{0.3250}$.
The \num{51.2}\,pp accuracy gain from the isolated GNN to SPECTRA-full
reflects the architectural design of the ensemble: the GraphSAGE component
(Module~G) operates on 16-dimensional node features that combine spectral
Laplacian embeddings with six behavioral statistics derived from each
address's transaction history.
When the GraphSAGE neighborhood aggregation is coupled with the VAE latent
codes (Module~B), the NMF behavioral profiles (Module~E), the Bayesian
trust scores (Module~D), and the Markov anomaly features (Module~F),
the node feature context is sufficiently enriched to discriminate
transaction types that are geometrically inseparable in the spectral
subspace alone.
The GNN module's low standalone AUC-ROC ($\num{0.4209}$, below chance for
a six-class problem) is diagnostic: the spectral embedding of the hub graph
does not by itself encode the fine-grained structural differences between,
say, Consolidation and Distribution transactions at the node level.
These differences only emerge when the ensemble integrates the richer
multi-modal feature context from the remaining six modules.

\textbf{F1-macro versus accuracy trade-off.}
EvolveGCN achieves $\text{F1}_{\text{macro}} = \num{0.5561}$, which
exceeds SPECTRA-full's $\text{F1}_{\text{macro}} = \num{0.3806}$.
This apparent contradiction is explained by the class distribution of the
node-level evaluation set.
The hub graph is heavily skewed toward exchange and mining-pool addresses
(Standard and Batch-Payment types); EvolveGCN's temporal graph convolution
correctly assigns more minority-class nodes on average, boosting macro-F1,
but simultaneously misclassifies a substantial fraction of the majority
class, suppressing accuracy and MCC.
SPECTRA-full, by contrast, achieves near-perfect identification of the
majority class at the cost of recall on rare categories, reflected in its
lower macro-F1 but substantially higher weighted-F1 ($\num{0.9134}$)
and MCC.
For forensic analysis, where false-positively flagging a legitimate exchange
address as anomalous carries operational cost, the high MCC and accuracy
profile of SPECTRA-full is preferable.

\textbf{AUC-ROC considerations.}
EvolveGCN's AUC-ROC of \num{0.8894} is notably higher than SPECTRA-full's
$\num{0.5545}$.
This gap reflects a calibration asymmetry: EvolveGCN outputs well-calibrated
class probabilities over its simpler two-layer temporal graph convolution,
yielding smooth ROC curves even when hard-threshold accuracy is moderate.
SPECTRA-full's ensemble pathway produces more confident (less probabilistic)
predictions, compressing the soft-score distribution and reducing area under
the ROC curve while raising hard-threshold metrics.
The AUC-ROC of the SPECTRA GNN module alone ($\num{0.4209}$) rising to
$\num{0.5545}$ after ensemble integration confirms that the auxiliary modules
contribute positive discriminative signal even to the probabilistic ranking.

% ─── VI-C  Temporally-Fair Comparison ────────────────────────────────────────

\subsection{Temporally-Fair Comparison}
\label{subsec:fair_results}

To test for temporal leakage, we remove the three absolute time identifiers
\texttt{block\_height}, \texttt{block\_time}, and \texttt{year} from all
tabular feature sets and re-evaluate every method.
Table~\ref{tab:fair_comparison} reports original and fair scores side-by-side.

\textbf{Key finding: temporal leakage is negligible for all methods.}
No method degrades substantially after removing temporal identifiers,
confirming that the performance differences observed in
Section~\ref{subsec:t1_results} and Section~\ref{subsec:t2_results}
are not artifacts of time-based shortcuts.

\textbf{XGBoost and BERT4ETH remain at ceiling.}
XGBoost retains $\text{Acc} = \num{1.0000}$ and $\text{F1}_{\text{macro}} = \num{1.0000}$
under the fair protocol; BERT4ETH drops to $\text{Acc} = \num{0.9998}$
($-\num{0.02}$\,pp), statistically indistinguishable from perfect.
This confirms the conclusion of Section~\ref{subsec:t1_results}: the
perfect scores arise from label-feature entanglement in the structural
fields, not from temporal shortcuts in \texttt{block\_height} or
\texttt{year}.
Removing temporal features does not break the algebraic correspondence
between feature values and labels.

\textbf{SPECTRA is unaffected by construction.}
SPECTRA operates on the address-level transaction graph and spectral
Laplacian embeddings; none of its seven modules consume \texttt{block\_height},
\texttt{block\_time}, or \texttt{year} as input features.
Its fair-protocol scores are therefore identical to the original scores,
satisfying the temporal fairness criterion by design rather than by
post-hoc ablation.

\textbf{EvolveGCN slightly improves under fair evaluation.}
EvolveGCN's accuracy rises from \num{0.5773} to \num{0.5844}
($+\num{0.71}$\,pp) and AUC-ROC from \num{0.8894} to \num{0.9063}
($+\num{1.69}$\,pp) after temporal feature removal.
The small improvement suggests a minor degree of temporal bias in the
original EvolveGCN evaluation: the temporal identifiers were providing
marginally confounding signal that, when removed, forces the model to
rely purely on graph topology---a modality it is better suited to exploit.
AA-GNN shows a similar mild improvement ($\text{Acc}: \num{0.4992} \to
\num{0.5115}$, $+\num{1.23}$\,pp), consistent with the same mechanism.

% ─── VI-D  Clustering Performance ────────────────────────────────────────────

\subsection{Clustering Performance}
\label{subsec:clustering_results}

Table~\ref{tab:clustering} reports clustering quality under four metrics:
Silhouette coefficient (Sil, $\uparrow$), Davies-Bouldin index
(DB, $\downarrow$), Calinski-Harabasz score (CH, $\uparrow$), and
Normalized Mutual Information with ground-truth labels (NMI, $\uparrow$).
Two evaluation conditions are shown: \emph{Original} (full feature set)
and \emph{Fair} (temporal identifiers removed).

\textbf{SPECTRA-HDBSCAN NMI is competitive with PlainHDBSCAN.}
SPECTRA-HDBSCAN achieves $\text{NMI} = \num{0.1118}$, closely matching
PlainHDBSCAN at \num{0.1176} ($-\num{0.58}$\,pp gap), while operating on
the 10-dimensional VAE latent space rather than the raw feature matrix.
The VAE compression preserves the cluster-relevant variance while
suppressing noise dimensions, yielding a representation that is
qualitatively richer even when NMI parity with the raw-space method holds.
No other clustering method simultaneously provides drift-aware temporal
segmentation, per-address trust certificates, and Markov behavioral
anomaly scores---capabilities that are orthogonal to the NMI metric
but critical for forensic interpretability.

\textbf{GraphAE achieves the highest silhouette score but lowest NMI.}
GraphAE records the best Silhouette ($\num{0.3419}$ original,
$\num{0.6903}$ fair) and the best Davies-Bouldin index in the fair condition
($\num{0.9406}$), indicating geometrically compact and well-separated
clusters in the embedding space.
However, its NMI of \num{0.0364} (original) and \num{0.0237} (fair) is the
lowest across all methods, demonstrating that geometric compactness does
not imply alignment with transaction-type ground truth.
GraphAE's graph autoencoder objective optimizes for structural
reconstruction of the adjacency matrix; the resulting embedding organizes
addresses by connectivity pattern rather than by transaction-type label,
producing tight topological communities that cross label boundaries.

\textbf{AE+KMeans achieves the best Calinski-Harabasz score.}
AE+KMeans attains $\text{CH} = \num{15213.5}$ (original) and
$\num{15451.3}$ (fair), the highest compact-cluster ratio among all methods,
confirming that its autoencoder latent space, when partitioned by K-Means,
yields globular clusters with high between-to-within scatter ratio.
The NMI of \num{0.0973}/\num{0.0905} is moderate, suggesting partial label
alignment.

\textbf{DBSCAN benefits substantially from temporal feature removal.}
DBSCAN's Silhouette improves from $-\num{0.0274}$ to $\num{0.3183}$ after
removing temporal identifiers ($+\num{0.346}$ absolute), and its
Davies-Bouldin index falls from \num{2.7686} to \num{0.9167}---the best
DB score across all fair-condition methods.
This dramatic shift indicates that temporal features introduced intra-cluster
heterogeneity that DBSCAN's density-connectivity criterion could not resolve;
without them, density-reachable neighborhoods become more homogeneous.

\textbf{Noise and structural properties of SPECTRA-HDBSCAN.}
The SPECTRA-HDBSCAN run identifies \num{135} clusters over the
\num{30000}-node hub graph, with \num{63390} points (\num{21.1}\%)
classified as noise.
The high noise fraction is consistent with the known heterogeneity of
Bitcoin hub addresses: a substantial fraction of exchange and mixer addresses
exhibit transaction patterns that are genuinely intermediate between
categories, and HDBSCAN correctly refuses to assign these ambiguous
nodes to any cluster rather than forcing them into artificial groupings.

% ─── VI-E  SPECTRA-Specific Diagnostic Outputs ───────────────────────────────

\subsection{SPECTRA-Specific Diagnostic Outputs}
\label{subsec:spectra_outputs}

The following outputs are produced exclusively by SPECTRA modules and
have no counterpart in any evaluated baseline.

\textbf{Wasserstein drift detection (Module~C).}
SPECTRA evaluates Wasserstein distance between consecutive temporal windows
of the VAE latent distribution.
Across all \num{29} evaluated windows, the drift detector flags
\emph{all \num{29}} as exceeding the alert threshold $\delta = 0.1$
(flagging rate: \num{100}\%).
Drift scores range from \num{0.179} to \num{50.745} with a mean of
$\bar{W} = \num{8.500}$, indicating that no two consecutive temporal
windows share a stationary distribution.
This finding is consistent with documented regime changes in the Bitcoin
ecosystem: the 2017 retail speculation cycle, the 2020--2021
DeFi/NFT expansion, and the 2022 market contraction each induced
fundamentally different transaction-graph topology and fee dynamics that
manifest as persistent concept drift in the latent space.
The magnitude variation (minimum \num{0.179} near stable periods,
maximum \num{50.745} during market dislocations) further suggests
that drift intensity correlates with macroeconomic Bitcoin events.
A static classifier trained on any single temporal window would therefore
exhibit significant distributional shift when deployed forward in time---
a risk that SPECTRA's continuous drift monitoring is designed to detect
and flag for model retraining.

\textbf{Bayesian trust scoring (Module~D).}
The Beta-distributed trust model assigns each hub address a posterior
trust score $\tau \sim \mathrm{Beta}(\alpha, \beta)$ conditioned on its
observed send/receive ratio, fee anomaly, and clustering membership.
The population-level mean trust is $\bar{\tau} = \num{0.634} \pm \num{0.105}$.
By the symmetry of the Beta distribution around $\tau = 0.5$, addresses
with $\tau < 0.5$ are classified as anomalous; under the observed mean and
variance, this corresponds to approximately \num{15.8}\% of the hub
population.
These $\approx \num{4740}$ anomalous addresses span exchanges with unusual
fee escalation patterns, structuring entities (addresses that fragment large
sums into just-below-threshold outputs), and probable mixer relays.
Crucially, the Bayesian framework produces calibrated uncertainty intervals
for each address rather than a hard binary label, enabling risk-tiered
investigation workflows in which analysts prioritize addresses by
probability mass below the anomaly threshold.

\textbf{Markov behavioral profiling (Module~F).}
SPECTRA fits a discrete-state Markov chain over the hub graph's temporal
transaction sequence, inferring \num{136} behavioral states.
The population mean anomaly score is $\bar{s}_{\text{Markov}} = \num{0.875}$.
A score near \num{1.0} indicates that the observed transition distribution
for that address is maximally entropic---indistinguishable from a
uniform random walk over the state space.
This entropy-maximizing behavior is the defining signature of high-throughput
exchanges and mixing services: their addresses transact with near-random
counterparties across the full hub-node population, producing transition
matrices that are near-uniform and thus assigned near-maximal anomaly score.
The finding that $\bar{s} = \num{0.875}$ for the \emph{hub} graph is
therefore expected and interpretable: by construction, the top-\num{30000}
degree nodes are precisely the exchanges, mining pools, and mixers that
exhibit this behavior.
Addresses falling substantially below the population mean (e.g.,
$s < 0.5$) exhibit structured, low-entropy transition patterns consistent
with purpose-specific wallets (payroll distributors, DCA bots, or
controlled consolidation services) and warrant closer forensic attention.

\textbf{VAE reconstruction anomaly threshold.}
The Variational Autoencoder (Module~B) reaches a best validation loss of
$\mathcal{L}_{\text{val}} = \num{0.2216}$ with a reconstruction-based
anomaly threshold of $\varepsilon_{\text{VAE}} = \num{0.1694}$.
Addresses whose reconstruction error exceeds $\varepsilon_{\text{VAE}}$
lie in low-density regions of the latent space and are flagged as
structurally anomalous independently of their Bayesian trust score or
Markov state assignment.
The three anomaly signals (VAE reconstruction, Bayesian trust, Markov
entropy) are deliberately non-redundant: VAE detects distributional
outliers, Bayesian trust captures behavioral irregularity relative to
peer addresses in the same cluster, and Markov scoring identifies
temporal transition anomalies.
An address triggering all three simultaneously represents the
highest-confidence anomaly designation in the SPECTRA pipeline.

% =============================================================================
%  Table I — Classification Performance (Tracks T1 and T2)
% =============================================================================

\begin{table*}[!t]
\centering
\caption{Classification performance on the held-out test set.
  Track~T1 (tabular) uses the \num{20}-feature set $\mathbf{X}^{\dagger}$.
  Track~T2 (graph-topology) uses the address graph $\mathcal{G}$ and
  spectral embeddings $\boldsymbol{\Phi}$ only.
  Bold entries indicate the best result in each column within each track.
  MCC: Matthews Correlation Coefficient; Inf: per-sample inference latency (ms).}
\label{tab:classification}
\begin{threeparttable}
\begin{tabular}{@{}l S[table-format=1.4] S[table-format=1.4] S[table-format=1.4]
                  S[table-format=1.4] S[table-format=1.4] S[table-format=1.4]
                  S[table-format=1.4] S[table-format=3.1] S[table-format=1.4]@{}}
\toprule
{Method} & {Acc} & {F1$_{\text{mac}}$} & {F1$_{\text{wt}}$}
         & {Prec} & {Rec} & {AUC-ROC} & {MCC}
         & {Train (s)} & {Inf (ms)} \\
\midrule
\multicolumn{10}{@{}l}{\textit{Track T1 — Tabular classifiers}} \\
\midrule
XGBoost\tnote{$\dagger$}
  & \textbf{1.0000} & \textbf{1.0000} & \textbf{1.0000}
  & \textbf{1.0000} & \textbf{1.0000} & \textbf{1.0000} & \textbf{1.0000}
  & 6.7 & \textbf{0.0013} \\
BERT4ETH\tnote{$\dagger$}
  & \textbf{1.0000} & \textbf{1.0000} & {---}
  & {---} & {---} & \textbf{1.0000} & \textbf{1.0000}
  & 229.5 & 0.0017 \\
\midrule
\multicolumn{10}{@{}l}{\textit{Track T2 — Graph-topology methods}} \\
\midrule
EvolveGCN~\cite{Pareja2020}
  & 0.5773 & \textbf{0.5561} & \textbf{0.5561}
  & \textbf{0.5815} & \textbf{0.5773} & \textbf{0.8894} & 0.5002
  & 2.18 & {---} \\
AA-GNN~\cite{AAGNN2021}
  & 0.4992 & 0.4800 & 0.4800
  & 0.5326 & 0.4992 & 0.8462 & 0.4101
  & 0.91 & {---} \\
SPECTRA (GNN only)
  & 0.4036 & 0.1549 & 0.4074
  & 0.2604 & 0.1931 & 0.4209 & 0.3250
  & 2.09 & {---} \\
SPECTRA-full
  & \textbf{0.9160} & 0.3806 & 0.9134
  & 0.3996 & 0.4009 & 0.5545 & \textbf{0.8397}
  & {---} & {---} \\
\bottomrule
\end{tabular}
\begin{tablenotes}[flushleft]\footnotesize
  \item[$\dagger$] Upper-bound oracle: labels are derived from the same
    structural features used as inputs; see Section~V-B.
    These scores do not represent generalizable predictive performance.
\end{tablenotes}
\end{threeparttable}
\end{table*}

% =============================================================================
%  Table II — Temporally-Fair Comparison
% =============================================================================

\begin{table}[!t]
\centering
\caption{Temporally-fair ablation.
  \emph{Original}: full feature set including \texttt{block\_height},
  \texttt{block\_time}, and \texttt{year}.
  \emph{Fair}: those three temporal identifiers removed.
  $\Delta$: fair minus original (pp).
  SPECTRA is unaffected by construction (temporal features not used).}
\label{tab:fair_comparison}
\begin{tabular}{@{}l
                S[table-format=1.4] S[table-format=1.4] S[table-format=+1.2]
                S[table-format=1.4] S[table-format=1.4] S[table-format=+1.2]@{}}
\toprule
& \multicolumn{3}{c}{Accuracy} & \multicolumn{3}{c}{F1$_{\text{macro}}$} \\
\cmidrule(lr){2-4}\cmidrule(lr){5-7}
{Method} & {Orig.} & {Fair} & {$\Delta$}
         & {Orig.} & {Fair} & {$\Delta$} \\
\midrule
XGBoost\tnote{$\dagger$}
  & 1.0000 & \textbf{1.0000} & {$\pm$0.00}
  & 1.0000 & \textbf{1.0000} & {$\pm$0.00} \\
BERT4ETH\tnote{$\dagger$}
  & 1.0000 & 0.9998 & {$-$0.02}
  & 1.0000 & 0.9998 & {$-$0.02} \\
EvolveGCN
  & 0.5773 & 0.5844 & {$+$0.71}
  & 0.5561 & 0.5692 & {$+$1.31} \\
AA-GNN
  & 0.4992 & 0.5115 & {$+$1.23}
  & 0.4800 & 0.4953 & {$+$1.53} \\
SPECTRA-full
  & \textbf{0.9160} & \textbf{0.9160} & {$\pm$0.00}
  & 0.3806 & 0.3806 & {$\pm$0.00} \\
\bottomrule
\end{tabular}
\end{table}

% =============================================================================
%  Table III — Clustering Performance
% =============================================================================

\begin{table*}[!t]
\centering
\caption{Clustering quality on the hub-graph node set ($|\mathcal{V}| = \num{30000}$).
  Sil: Silhouette coefficient ($\uparrow$~better);
  DB: Davies-Bouldin index ($\downarrow$~better);
  CH: Calinski-Harabasz score ($\uparrow$~better);
  NMI: Normalized Mutual Information with ground-truth labels ($\uparrow$~better).
  \emph{Orig.}: original feature set; \emph{Fair}: temporal identifiers removed.
  Bold entries indicate the best result in each column across all methods
  within each evaluation condition.}
\label{tab:clustering}
\begin{tabular}{@{}l
                S[table-format=+1.4] S[table-format=1.4]
                S[table-format=5.1]  S[table-format=1.4]
                S[table-format=+1.4] S[table-format=1.4]
                S[table-format=5.1]  S[table-format=1.4]@{}}
\toprule
& \multicolumn{4}{c}{\textit{Original}} & \multicolumn{4}{c}{\textit{Fair}} \\
\cmidrule(lr){2-5}\cmidrule(lr){6-9}
{Method}
  & {Sil $\uparrow$} & {DB $\downarrow$} & {CH $\uparrow$} & {NMI $\uparrow$}
  & {Sil $\uparrow$} & {DB $\downarrow$} & {CH $\uparrow$} & {NMI $\uparrow$} \\
\midrule
K-Means~\cite{Lloyd1982}
  & 0.1675 & 1.4900 & 7050.8 & 0.0430
  & 0.1727 & 1.3214 & 7457.7 & 0.0307 \\
DBSCAN~\cite{Ester1996}
  & -0.0274 & 2.7686 & 739.8 & 0.0990
  & 0.3183 & \textbf{0.9167} & 3290.5 & 0.1079 \\
Isolation Forest~\cite{Liu2008}
  & {---} & {---} & {---} & 0.0217
  & {---} & {---} & {---} & {---} \\
AE+KMeans
  & 0.2829 & 1.1032 & \textbf{15213.5} & 0.0973
  & 0.3249 & 1.2796 & \textbf{15451.3} & 0.0905 \\
PlainHDBSCAN~\cite{McInnes2017}
  & 0.2013 & 2.4074 & 2612.8 & \textbf{0.1176}
  & \textbf{0.3778} & 1.0369 & 6603.9 & \textbf{0.1227} \\
GraphAE~\cite{Kipf2016VAE}
  & \textbf{0.3419} & \textbf{1.1024} & 16544.9 & 0.0364
  & \textbf{0.6903} & 0.9406 & 18013.6 & 0.0237 \\
SPECTRA-HDBSCAN
  & 0.0177 & 3.5805 & 5595.8 & 0.1118
  & 0.0177 & 3.5805 & 5595.8 & 0.1118 \\
\bottomrule
\end{tabular}
\end{table*}

% =============================================================================
%  Summary paragraph
% =============================================================================

\subsection{Summary of Findings}
\label{subsec:summary}

The results across all four evaluation axes converge on three principal
conclusions.
First, tabular classifiers (Track~T1) achieve ceiling performance due to
label-feature entanglement rather than genuine generalization; this is
confirmed by the temporally-fair ablation and the absence of performance
drop when time identifiers are removed.
Second, among all methods operating under the methodologically sound
graph-topology regime (Track~T2), SPECTRA-full achieves the highest accuracy
(\num{0.916}) and MCC (\num{0.840}), demonstrating that multi-modal feature
fusion across seven modules substantially outperforms single-objective
graph neural networks on this benchmark.
Third, SPECTRA provides unique forensic value beyond any classification or
clustering metric: continuous Wasserstein drift monitoring (100\% alert rate
over \num{29} windows), calibrated Bayesian trust certificates
($\bar{\tau} = \num{0.634} \pm \num{0.105}$, $\approx$\,15.8\% anomalous),
and Markov behavioral profiling (\num{136} states, $\bar{s} = \num{0.875}$)
collectively constitute an interpretable, multi-layered forensic intelligence
layer that no single-objective baseline provides.
